\documentclass[11pt,a4paper]{article}
\usepackage[utf8]{inputenc}
\usepackage[T1]{fontenc}
\usepackage[spanish,es-minimal,es-noshorthands]{babel}
\AtBeginDocument{\renewcommand{\sin}{\operatorname{sen}}\renewcommand{\arcsin}{\operatorname{arcsen}}}
\usepackage{lmodern}
\usepackage[sfdefault,lining]{FiraSans}
\usepackage[a4paper,top=2.3cm,bottom=2.2cm,left=1.8cm,right=1.8cm,headheight=1.3cm,headsep=0.4cm,footskip=1cm]{geometry}
\usepackage{amsmath,amssymb}
\usepackage{icomma}
\usepackage{graphicx}
\usepackage[table]{xcolor}
\usepackage{tikz}
\usepackage{fontawesome5}
\usepackage{ccicons}
\usepackage{fancyhdr}
\usepackage{lastpage}
\usepackage{tasks}
\usepackage{needspace}
\usepackage[most]{tcolorbox}
\usepackage{array}
\usepackage[hidelinks]{hyperref}
\AtBeginDocument{\hypersetup{pdftitle={Soluciones de la ficha de la unidad 7: Geometría básica del plano. 1º ESO},pdfsubject={Matemáticas 1º ESO},pdfkeywords={punto, recta, semirrecta, segmento, rectas perpendiculares, mediatriz, ángulos, bisectriz, sistema sexagesimal, complementarios, suplementarios, polígonos, diagonales, triángulos, cuadriláteros, polígonos regulares}}}

\definecolor{navy}{HTML}{1B2748}
\definecolor{accent}{HTML}{E8603A}
\definecolor{teal}{HTML}{1F9E90}
\definecolor{soft}{HTML}{EEF1F8}
\definecolor{muted}{HTML}{5D6685}

\setlength{\parindent}{0pt}
\setlength{\parskip}{3pt}
\settasks{label=\alph*), label-format=\bfseries\color{navy}, label-width=1.3em, item-indent=1.8em, column-sep=0.6em, after-item-skip=5pt, before-skip=3pt, after-skip=4pt}

% ---------- capçalera i peu ----------
\newcommand{\llicencia}{\href{https://creativecommons.org/licenses/by-nc-sa/4.0/deed.es}{\ccbyncsa}\ \ Licencia Creative Commons Atribución-NoComercial-CompartirIgual 4.0 Internacional (CC BY-NC-SA 4.0)}
\pagestyle{fancy}
\fancyhf{}
\fancyhead[L]{\raisebox{-0.35\height}{\includegraphics[height=0.95cm]{logo-gva.png}}}
\fancyhead[C]{\small\color{muted}Soluciones de la ficha · Unidad 7 · 1\textsuperscript{o} ESO}
\fancyhead[R]{\raisebox{-0.35\height}{\includegraphics[height=0.95cm]{logo-vermellar.png}}}
\fancyfoot[L]{\scriptsize\color{muted}\llicencia}
\fancyfoot[R]{\small\color{muted}\thepage}
\renewcommand{\headrulewidth}{0.4pt}
\renewcommand{\headrule}{\color{muted!40}\hrule height \headrulewidth}
\fancypagestyle{primera}{\fancyhf{}\renewcommand{\headrulewidth}{0pt}
  \fancyfoot[L]{\scriptsize\color{muted}\llicencia}
  \fancyfoot[R]{\small\color{muted}\thepage}}

% ---------- nivells i exercicis ----------
\newcommand{\estrelles}[1]{{\color{accent}\ifcase#1\or\faStar\faStar[regular]\faStar[regular]\or\faStar\faStar\faStar[regular]\or\faStar\faStar\faStar\fi}}
\newcounter{exer}
\newcommand{\ex}[1]{\par\needspace{4\baselineskip}\medskip\stepcounter{exer}\noindent{\color{navy}\bfseries\theexer.}\ \ {\small\estrelles{#1}}\ \ }

\newcommand{\seccio}[2]{\par\needspace{10\baselineskip}\bigskip
  \begin{tcolorbox}[colback=navy,colframe=navy,coltext=white,arc=2mm,boxrule=0pt,left=3mm,right=3mm,top=1.2mm,bottom=1.2mm]
  {\large\bfseries #1\quad #2}
  \end{tcolorbox}\vspace{-2pt}}

\newcolumntype{C}[1]{>{\centering\arraybackslash}p{#1}}

% quadre de recordatori (només a la fitxa d'exercicis)
\newcommand{\recorda}[1]{\par\smallskip\begin{tcolorbox}[colback=accent!7,colframe=accent!55,arc=2mm,boxrule=0.5pt,left=2.5mm,right=2.5mm,top=1mm,bottom=1mm]\small{\color{accent}\faLightbulb[regular]}\ \textbf{Recuerda:}\ #1\end{tcolorbox}}
\newcommand{\res}[1]{\boxed{#1}}\newcommand{\rest}[1]{\fbox{#1}}

\def\fmneed{14}
% =====================================================================
% Capçalera FentMates de les fitxes: un color per a cada curs.
% Els colors són els de la pàgina inicial (LEVELS a assets/js/materials.js);
% el logo ve de _eines/logo/pdf/ (python3 _eines/logo/genera.py).
%
% Ús, al final del preàmbul de la unitat:
%   % =====================================================================
% Capçalera FentMates de les fitxes: un color per a cada curs.
% Els colors són els de la pàgina inicial (LEVELS a assets/js/materials.js);
% el logo ve de _eines/logo/pdf/ (python3 _eines/logo/genera.py).
%
% Ús, al final del preàmbul de la unitat:
%   % =====================================================================
% Capçalera FentMates de les fitxes: un color per a cada curs.
% Els colors són els de la pàgina inicial (LEVELS a assets/js/materials.js);
% el logo ve de _eines/logo/pdf/ (python3 _eines/logo/genera.py).
%
% Ús, al final del preàmbul de la unitat:
%   \input{../comu/fentmates.tex}
%   \fmcurs{1eso}
% i, al començament del document:
%   \capcalera{Matemàtiques}{1r ESO}{Fitxa d'exercicis · Unitat 2}{Divisibilitat}
% Opcional abans de \input: \def\fmneed{14} (espai mínim abans de cada secció).
% =====================================================================
\makeatletter
\def\fm@logodir{../../logo/pdf}
\providecommand{\fmneed}{10}
\def\fm@cursnom{}

% clau del curs, color i color del text damunt del color
\newcommand{\fm@curs}[3]{\@namedef{fm@color@#1}{#2}\@namedef{fm@text@#1}{#3}}
\fm@curs{1eso}{E6B84A}{navy}
\fm@curs{2eso}{3B82F6}{white}
\fm@curs{3eso}{FF7A50}{white}
\fm@curs{4eso-a}{A78BFA}{white}
\fm@curs{4eso-b}{6D28D9}{white}
\fm@curs{1batx-mi}{9B1C1C}{white}
\fm@curs{1batx-ccss}{EF6B6B}{white}
\fm@curs{2batx-mii}{1E6B47}{white}
\fm@curs{2batx-ccss}{2BB3A3}{white}

\newcommand{\fmcurs}[1]{%
  \@ifundefined{fm@color@#1}{\PackageError{fentmates}{Curs desconegut: #1}{}}{}%
  \edef\fm@tmp{\noexpand\definecolor{curs}{HTML}{\@nameuse{fm@color@#1}}%
               \noexpand\colorlet{curstext}{\@nameuse{fm@text@#1}}}\fm@tmp
  \colorlet{cursfosc}{curs!72!black}%
  \def\fm@logo{\fm@logodir/vertical-#1.pdf}%
  \def\fm@marca{\fm@logodir/marca-#1.pdf}}

% ---------- capçalera de la primera pàgina ----------
% #1 matèria (petit), #2 curs, #3 tipus de document i unitat, #4 títol
\newcommand{\capcalera}[4]{%
  \gdef\fm@cursnom{#2}%
  \par\noindent
  \begin{tcolorbox}[colback=curs,colframe=curs,boxrule=0pt,arc=2.2mm,width=3.5cm,height=2.15cm,
      valign=center,left=3mm,right=2mm,top=1mm,bottom=1mm,nobeforeafter,box align=center]
    \color{curstext}{\footnotesize #1}\par\vspace{1pt}{\fontsize{19}{21}\selectfont\bfseries #2\par}%
  \end{tcolorbox}\hspace{5mm}%
  \begin{minipage}[c]{\dimexpr\linewidth-3.5cm-5mm-2.9cm\relax}
    {\color{cursfosc}\small\bfseries\MakeUppercase{#3}}\par\vspace{2pt}
    {\color{navy}\fontsize{20}{23}\selectfont\bfseries #4\par}
  \end{minipage}\hfill
  \begin{minipage}[c]{2.6cm}\raggedleft\includegraphics[width=2.6cm]{\fm@logo}\end{minipage}\par
  \vspace{2pt}}

% ---------- barres de secció del color del curs ----------
\renewcommand{\seccio}[2]{\par\needspace{\fmneed\baselineskip}\bigskip
  \begin{tcolorbox}[colback=curs,colframe=curs,coltext=curstext,arc=2mm,boxrule=0pt,left=3mm,right=3mm,top=1.2mm,bottom=1.2mm]
  {\large\bfseries #1\quad #2}
  \end{tcolorbox}\vspace{-2pt}}

% ---------- estrelles de dificultat del color del curs (a Batxillerat no se'n posen) ----------
\@ifundefined{estrelles}{}{%
  \renewcommand{\estrelles}[1]{{\color{curs}\ifcase#1\or\faStar\faStar[regular]\faStar[regular]\or\faStar\faStar\faStar[regular]\or\faStar\faStar\faStar\fi}}}

% ---------- capçalera i peu de les pàgines ----------
% primera pàgina: logos institucionals al peu
\fancypagestyle{primera}{\fancyhf{}\renewcommand{\headrulewidth}{0pt}
  \fancyfoot[L]{\raisebox{-0.32\height}{\includegraphics[height=0.78cm]{logo-gva.png}}\hspace{5mm}%
    \raisebox{-0.32\height}{\includegraphics[height=0.78cm]{logo-vermellar.png}}}
  \fancyfoot[R]{\scriptsize\color{muted}\href{https://creativecommons.org/licenses/by-nc-sa/4.0/deed.ca}{\ccbyncsa}\ \ CC BY-NC-SA 4.0 · fentmates.org\qquad\small\thepage}}
% la resta: el símbol a l'esquerra i el curs a la dreta, amb el títol al centre com abans
\fancyhead[L]{\raisebox{-0.3\height}{\includegraphics[height=0.62cm]{\fm@marca}}}
\fancyhead[R]{\small\bfseries\color{cursfosc}\fm@cursnom}
\renewcommand{\headrule}{\color{curs!45}\hrule height \headrulewidth}
\makeatother

%   \fmcurs{1eso}
% i, al començament del document:
%   \capcalera{Matemàtiques}{1r ESO}{Fitxa d'exercicis · Unitat 2}{Divisibilitat}
% Opcional abans de \input: \def\fmneed{14} (espai mínim abans de cada secció).
% =====================================================================
\makeatletter
\def\fm@logodir{../../logo/pdf}
\providecommand{\fmneed}{10}
\def\fm@cursnom{}

% clau del curs, color i color del text damunt del color
\newcommand{\fm@curs}[3]{\@namedef{fm@color@#1}{#2}\@namedef{fm@text@#1}{#3}}
\fm@curs{1eso}{E6B84A}{navy}
\fm@curs{2eso}{3B82F6}{white}
\fm@curs{3eso}{FF7A50}{white}
\fm@curs{4eso-a}{A78BFA}{white}
\fm@curs{4eso-b}{6D28D9}{white}
\fm@curs{1batx-mi}{9B1C1C}{white}
\fm@curs{1batx-ccss}{EF6B6B}{white}
\fm@curs{2batx-mii}{1E6B47}{white}
\fm@curs{2batx-ccss}{2BB3A3}{white}

\newcommand{\fmcurs}[1]{%
  \@ifundefined{fm@color@#1}{\PackageError{fentmates}{Curs desconegut: #1}{}}{}%
  \edef\fm@tmp{\noexpand\definecolor{curs}{HTML}{\@nameuse{fm@color@#1}}%
               \noexpand\colorlet{curstext}{\@nameuse{fm@text@#1}}}\fm@tmp
  \colorlet{cursfosc}{curs!72!black}%
  \def\fm@logo{\fm@logodir/vertical-#1.pdf}%
  \def\fm@marca{\fm@logodir/marca-#1.pdf}}

% ---------- capçalera de la primera pàgina ----------
% #1 matèria (petit), #2 curs, #3 tipus de document i unitat, #4 títol
\newcommand{\capcalera}[4]{%
  \gdef\fm@cursnom{#2}%
  \par\noindent
  \begin{tcolorbox}[colback=curs,colframe=curs,boxrule=0pt,arc=2.2mm,width=3.5cm,height=2.15cm,
      valign=center,left=3mm,right=2mm,top=1mm,bottom=1mm,nobeforeafter,box align=center]
    \color{curstext}{\footnotesize #1}\par\vspace{1pt}{\fontsize{19}{21}\selectfont\bfseries #2\par}%
  \end{tcolorbox}\hspace{5mm}%
  \begin{minipage}[c]{\dimexpr\linewidth-3.5cm-5mm-2.9cm\relax}
    {\color{cursfosc}\small\bfseries\MakeUppercase{#3}}\par\vspace{2pt}
    {\color{navy}\fontsize{20}{23}\selectfont\bfseries #4\par}
  \end{minipage}\hfill
  \begin{minipage}[c]{2.6cm}\raggedleft\includegraphics[width=2.6cm]{\fm@logo}\end{minipage}\par
  \vspace{2pt}}

% ---------- barres de secció del color del curs ----------
\renewcommand{\seccio}[2]{\par\needspace{\fmneed\baselineskip}\bigskip
  \begin{tcolorbox}[colback=curs,colframe=curs,coltext=curstext,arc=2mm,boxrule=0pt,left=3mm,right=3mm,top=1.2mm,bottom=1.2mm]
  {\large\bfseries #1\quad #2}
  \end{tcolorbox}\vspace{-2pt}}

% ---------- estrelles de dificultat del color del curs (a Batxillerat no se'n posen) ----------
\@ifundefined{estrelles}{}{%
  \renewcommand{\estrelles}[1]{{\color{curs}\ifcase#1\or\faStar\faStar[regular]\faStar[regular]\or\faStar\faStar\faStar[regular]\or\faStar\faStar\faStar\fi}}}

% ---------- capçalera i peu de les pàgines ----------
% primera pàgina: logos institucionals al peu
\fancypagestyle{primera}{\fancyhf{}\renewcommand{\headrulewidth}{0pt}
  \fancyfoot[L]{\raisebox{-0.32\height}{\includegraphics[height=0.78cm]{logo-gva.png}}\hspace{5mm}%
    \raisebox{-0.32\height}{\includegraphics[height=0.78cm]{logo-vermellar.png}}}
  \fancyfoot[R]{\scriptsize\color{muted}\href{https://creativecommons.org/licenses/by-nc-sa/4.0/deed.ca}{\ccbyncsa}\ \ CC BY-NC-SA 4.0 · fentmates.org\qquad\small\thepage}}
% la resta: el símbol a l'esquerra i el curs a la dreta, amb el títol al centre com abans
\fancyhead[L]{\raisebox{-0.3\height}{\includegraphics[height=0.62cm]{\fm@marca}}}
\fancyhead[R]{\small\bfseries\color{cursfosc}\fm@cursnom}
\renewcommand{\headrule}{\color{curs!45}\hrule height \headrulewidth}
\makeatother

%   \fmcurs{1eso}
% i, al començament del document:
%   \capcalera{Matemàtiques}{1r ESO}{Fitxa d'exercicis · Unitat 2}{Divisibilitat}
% Opcional abans de \input: \def\fmneed{14} (espai mínim abans de cada secció).
% =====================================================================
\makeatletter
\def\fm@logodir{../../logo/pdf}
\providecommand{\fmneed}{10}
\def\fm@cursnom{}

% clau del curs, color i color del text damunt del color
\newcommand{\fm@curs}[3]{\@namedef{fm@color@#1}{#2}\@namedef{fm@text@#1}{#3}}
\fm@curs{1eso}{E6B84A}{navy}
\fm@curs{2eso}{3B82F6}{white}
\fm@curs{3eso}{FF7A50}{white}
\fm@curs{4eso-a}{A78BFA}{white}
\fm@curs{4eso-b}{6D28D9}{white}
\fm@curs{1batx-mi}{9B1C1C}{white}
\fm@curs{1batx-ccss}{EF6B6B}{white}
\fm@curs{2batx-mii}{1E6B47}{white}
\fm@curs{2batx-ccss}{2BB3A3}{white}

\newcommand{\fmcurs}[1]{%
  \@ifundefined{fm@color@#1}{\PackageError{fentmates}{Curs desconegut: #1}{}}{}%
  \edef\fm@tmp{\noexpand\definecolor{curs}{HTML}{\@nameuse{fm@color@#1}}%
               \noexpand\colorlet{curstext}{\@nameuse{fm@text@#1}}}\fm@tmp
  \colorlet{cursfosc}{curs!72!black}%
  \def\fm@logo{\fm@logodir/vertical-#1.pdf}%
  \def\fm@marca{\fm@logodir/marca-#1.pdf}}

% ---------- capçalera de la primera pàgina ----------
% #1 matèria (petit), #2 curs, #3 tipus de document i unitat, #4 títol
\newcommand{\capcalera}[4]{%
  \gdef\fm@cursnom{#2}%
  \par\noindent
  \begin{tcolorbox}[colback=curs,colframe=curs,boxrule=0pt,arc=2.2mm,width=3.5cm,height=2.15cm,
      valign=center,left=3mm,right=2mm,top=1mm,bottom=1mm,nobeforeafter,box align=center]
    \color{curstext}{\footnotesize #1}\par\vspace{1pt}{\fontsize{19}{21}\selectfont\bfseries #2\par}%
  \end{tcolorbox}\hspace{5mm}%
  \begin{minipage}[c]{\dimexpr\linewidth-3.5cm-5mm-2.9cm\relax}
    {\color{cursfosc}\small\bfseries\MakeUppercase{#3}}\par\vspace{2pt}
    {\color{navy}\fontsize{20}{23}\selectfont\bfseries #4\par}
  \end{minipage}\hfill
  \begin{minipage}[c]{2.6cm}\raggedleft\includegraphics[width=2.6cm]{\fm@logo}\end{minipage}\par
  \vspace{2pt}}

% ---------- barres de secció del color del curs ----------
\renewcommand{\seccio}[2]{\par\needspace{\fmneed\baselineskip}\bigskip
  \begin{tcolorbox}[colback=curs,colframe=curs,coltext=curstext,arc=2mm,boxrule=0pt,left=3mm,right=3mm,top=1.2mm,bottom=1.2mm]
  {\large\bfseries #1\quad #2}
  \end{tcolorbox}\vspace{-2pt}}

% ---------- estrelles de dificultat del color del curs (a Batxillerat no se'n posen) ----------
\@ifundefined{estrelles}{}{%
  \renewcommand{\estrelles}[1]{{\color{curs}\ifcase#1\or\faStar\faStar[regular]\faStar[regular]\or\faStar\faStar\faStar[regular]\or\faStar\faStar\faStar\fi}}}

% ---------- capçalera i peu de les pàgines ----------
% primera pàgina: logos institucionals al peu
\fancypagestyle{primera}{\fancyhf{}\renewcommand{\headrulewidth}{0pt}
  \fancyfoot[L]{\raisebox{-0.32\height}{\includegraphics[height=0.78cm]{logo-gva.png}}\hspace{5mm}%
    \raisebox{-0.32\height}{\includegraphics[height=0.78cm]{logo-vermellar.png}}}
  \fancyfoot[R]{\scriptsize\color{muted}\href{https://creativecommons.org/licenses/by-nc-sa/4.0/deed.ca}{\ccbyncsa}\ \ CC BY-NC-SA 4.0 · fentmates.org\qquad\small\thepage}}
% la resta: el símbol a l'esquerra i el curs a la dreta, amb el títol al centre com abans
\fancyhead[L]{\raisebox{-0.3\height}{\includegraphics[height=0.62cm]{\fm@marca}}}
\fancyhead[R]{\small\bfseries\color{cursfosc}\fm@cursnom}
\renewcommand{\headrule}{\color{curs!45}\hrule height \headrulewidth}
\makeatother

\fmcurs{1eso}

\begin{document}
\thispagestyle{primera}

\capcalera{Matemáticas}{1\textsuperscript{o} ESO}{Soluciones de la ficha · Unidad 7}{Geometría básica del plano}
\vspace{8pt}

\begin{tcolorbox}[colback=soft,colframe=soft,arc=2mm,left=3mm,right=3mm,top=1.5mm,bottom=1.5mm]
\small Intenta resolver cada ejercicio antes de mirar su solución. Si no te sale el mismo resultado, repasa los pasos o consulta los apuntes interactivos de la unidad: encontrarás una calculadora de ángulos en el sistema sexagesimal y herramientas para explorar rectas, ángulos, triángulos y polígonos.
\end{tcolorbox}

% =====================================================
\seccio{1}{Puntos, rectas, semirrectas y segmentos}

\ex{1} \begin{tasks}(1)
  \task $A$, $B$ y $C$ son de la recta $r$ (están \rest{alineados}); $D$ \rest{no} está en ella.
  \task \rest{$\overline{AB}$, $\overline{AC}$ y $\overline{BC}$} (3 segmentos).
  \task \rest{Dos}: una hacia $A$ y otra hacia $C$.
\end{tasks}

\ex{1} \begin{tasks}(3)
  \task \rest{Recta} (dos flechas)
  \task \rest{Segmento} (dos extremos)
  \task \rest{Semirrecta} (una flecha)
\end{tasks}

\ex{1} \begin{tasks}(1)
  \task $7:2=\res{3,5\text{ cm}}$
  \task $2\cdot4,2=\res{8,4\text{ cm}}$
\end{tasks}

\ex{2} Con 4 puntos: $3+2+1=\res{6}$ segmentos. Con 6 puntos: $5+4+3+2+1=\res{15}$ segmentos.

\ex{2} En el primer caso, $\overline{AC}=5+3=\res{8\text{ cm}}$. En el segundo, $\overline{AC}=5-3=\res{2\text{ cm}}$.

\ex{3} De cada punto salen tantas rectas como puntos quedan, y cada recta se cuenta dos veces: \ $\frac{4\cdot3}{2}=\res{6}$, \ $\frac{5\cdot4}{2}=\res{10}$, \ $\frac{7\cdot6}{2}=\res{21}$.


% =====================================================
\seccio{2}{Posiciones relativas de rectas. Mediatriz}

\ex{1} \begin{tasks}(4)
  \task \rest{paralelas}
  \task \rest{perpendiculares}
  \task \rest{perpendiculares}
  \task \rest{secantes}
\end{tasks}

\ex{1} \begin{tasks}(2)
  \task \rest{Paralelas}
  \task \rest{Perpendiculares}
  \task \rest{Perpendiculares} (ángulo de $90^\circ$)
  \task \rest{Secantes} (ángulo de $60^\circ$)
\end{tasks}

\ex{1} Con el compás abierto más de 3\,cm, traza arcos desde $A$ y desde $B$ que se corten por encima y por debajo del segmento; la recta que une los dos puntos de corte es la mediatriz. Corta el segmento en el punto medio, a \rest{3\,cm} de $A$, y es perpendicular al segmento.

\ex{2} \begin{tasks}(1)
  \task \rest{Cierto}: se cortan en un punto.
  \task \rest{Falso}: solo si forman ángulos rectos.
  \task \rest{Cierto}.
  \task \rest{Falso}: $r$ y $t$ son paralelas.
\end{tasks}

\ex{2} Los puntos de la mediatriz equidistan de los extremos: $\overline{PB}=\res{5\text{ cm}}$. \ $\overline{AM}=8:2=\res{4\text{ cm}}$.

\ex{3} Los puntos que equidistan de $A$ y $B$ son los de la \rest{mediatriz de $\overline{AB}$}. La fuente tiene que estar en el \rest{punto donde la mediatriz corta el camino}.


% =====================================================
\seccio{3}{Ángulos: elementos, medida y clasificación}

\ex{1} \begin{tasks}(4)
  \task \rest{agudo}
  \task \rest{recto}
  \task \rest{obtuso}
  \task \rest{llano}
  \task \rest{cóncavo}
  \task \rest{completo}
  \task \rest{obtuso}
  \task \rest{agudo}
\end{tasks}

\ex{1} \begin{tasks}(4)
  \task \rest{agudo} ($40^\circ$)
  \task \rest{recto} ($90^\circ$)
  \task \rest{obtuso} ($140^\circ$)
  \task \rest{cóncavo} ($230^\circ$)
\end{tasks}

\ex{1} Entre dos números de la esfera hay $360^\circ:12=30^\circ$.\begin{tasks}(6)
  \task $\res{90^\circ}$
  \task $\res{120^\circ}$
  \task $\res{150^\circ}$
  \task $\res{180^\circ}$
  \task $\res{90^\circ}$
  \task $\res{30^\circ}$
\end{tasks}

\ex{2} Cada minuto gira $360^\circ:60=6^\circ$.\begin{tasks}(4)
  \task $\res{60^\circ}$
  \task $\res{150^\circ}$
  \task $\res{360^\circ}$
  \task $45\cdot6=\res{270^\circ}$
\end{tasks}

\ex{2} \begin{tasks}(1)
  \task $136^\circ:2=\res{68^\circ}$ cada uno.
  \task $2\cdot43^\circ=\res{86^\circ}$ (agudo).
\end{tasks}

\ex{2} Entre los dos forman una vuelta completa: $360^\circ-125^\circ=\res{235^\circ}$.

\ex{3} \begin{tasks}(1)
  \task $360^\circ:5=\res{72^\circ}$; \ $360^\circ:8=\res{45^\circ}$.
  \task $180:15=\res{12}$ veces.
\end{tasks}


% =====================================================
\seccio{4}{Medida de ángulos: el sistema sexagesimal}

\ex{1} \begin{tasks}(3)
  \task $\res{7200''}$
  \task $\res{900''}$
  \task $\res{4800''}$
  \task $\res{11140''}$
  \task $\res{36610''}$
  \task $\res{2730''}$
\end{tasks}

\ex{1} \begin{tasks}(4)
  \task $\res{300'}$
  \task $\res{150'}$
  \task $\res{765'}$
  \task $\res{10\,800'}$
\end{tasks}

\ex{1} Dividimos entre 60 tantas veces como haga falta: por ejemplo, $4000'':60=66'$ y sobran $40''$; $66':60=1^\circ$ y sobran $6'$.\begin{tasks}(3)
  \task $\res{2^\circ\,30'}$
  \task $\res{1^\circ\,6'\,40''}$
  \task $\res{2^\circ\,3'\,5''}$
  \task $\res{2^\circ\,46'\,40''}$
  \task $\res{16^\circ\,40'}$
  \task $\res{24^\circ}$
\end{tasks}

\ex{2} La parte decimal de los grados se multiplica por 60: $0,5^\circ=30'$, \ $0,25^\circ=15'$, \ $0,75^\circ=45'$, \ $0,1^\circ=6'$.\begin{tasks}(4)
  \task $\res{12^\circ\,30'}$
  \task $\res{30^\circ\,15'}$
  \task $\res{8^\circ\,45'}$
  \task $\res{6'}$
\end{tasks}

\ex{2} \begin{tasks}(2)
  \task $23^\circ\,45'+18^\circ\,30'=\res{42^\circ\,15'}$
  \task $45^\circ\,36'\,50''+27^\circ\,48'\,25''=\res{73^\circ\,25'\,15''}$
  \task $120^\circ\,15'\,40''+59^\circ\,44'\,20''=\res{180^\circ}$
\end{tasks}

\ex{2} Cuando no se puede restar, convertimos $1^\circ$ en $60'$ o $1'$ en $60''$.\begin{tasks}(2)
  \task $75^\circ\,40'-32^\circ\,15'=\res{43^\circ\,25'}$
  \task $60^\circ\,20'\,15''-25^\circ\,45'\,30''=\res{34^\circ\,34'\,45''}$
  \task $180^\circ-97^\circ\,34'\,18''=\res{82^\circ\,25'\,42''}$
\end{tasks}

\ex{2} \begin{tasks}(2)
  \task $(12^\circ\,20')\cdot3=\res{37^\circ}$
  \task $(25^\circ\,15'\,40'')\cdot4=\res{101^\circ\,2'\,40''}$
  \task $(8^\circ\,45'\,50'')\cdot6=\res{52^\circ\,35'}$
\end{tasks}

\ex{3} \begin{tasks}(2)
  \task $(75^\circ\,36'):4=\res{18^\circ\,54'}$
  \task $(100^\circ\,20'\,30''):5=\res{20^\circ\,4'\,6''}$
  \task $(47^\circ\,13'):3=\res{15^\circ\,44'\,20''}$
\end{tasks}Comprobación del primero: $(18^\circ\,54')\cdot4=72^\circ\,216'=75^\circ\,36'$ \checkmark

\ex{3} En 1 hora: $360^\circ:24=\res{15^\circ}$. En 1 minuto: $15^\circ:60=900':60=\res{15'}$. En 1 segundo: $15':60=900'':60=\res{15''}$.


% =====================================================
\seccio{5}{Relaciones entre ángulos}

\ex{1} \begin{tasks}(4)
  \task $\res{60^\circ}$
  \task $\res{43^\circ}$
  \task $\res{1^\circ}$
  \task $\res{74^\circ\,30'}$
\end{tasks}

\ex{1} \begin{tasks}(4)
  \task $\res{135^\circ}$
  \task $\res{60^\circ}$
  \task $\res{1^\circ}$
  \task $\res{87^\circ\,44'\,20''}$
\end{tasks}

\ex{1} $a$ y $c$ son adyacentes al ángulo de $35^\circ$: $a=c=180^\circ-35^\circ=\res{145^\circ}$. \ $b$ es opuesto por el vértice al ángulo de $35^\circ$: $b=\res{35^\circ}$.

\ex{2} El ángulo 3 es opuesto por el vértice al 1 y los ángulos 5 y 7 son correspondientes al 1 y al 3: \ $\hat3=\hat5=\hat7=\res{58^\circ}$. \ Los ángulos 2, 4, 6 y 8 son suplementarios del 1: \ $180^\circ-58^\circ=\res{122^\circ}$.

\ex{2} \begin{tasks}(1)
  \task \rest{Falso}: dos ángulos agudos suman menos de $180^\circ$.
  \task \rest{Cierto}: $180^\circ$ menos un ángulo agudo da más de $90^\circ$.
  \task \rest{Cierto}.
  \task \rest{Falso}: ya pasa de $90^\circ$.
\end{tasks}

\ex{2} $89^\circ\,59'\,60''-37^\circ\,25'\,48''=\res{52^\circ\,34'\,12''}$ \qquad $179^\circ\,59'\,60''-37^\circ\,25'\,48''=\res{142^\circ\,34'\,12''}$

\ex{3} $x+(x+20)=90\ \Rightarrow\ 2x=70\ \Rightarrow\ x=35$. \ Los ángulos miden \rest{$35^\circ$ y $55^\circ$}.

\ex{3} $x+4x=180\ \Rightarrow\ 5x=180\ \Rightarrow\ x=\res{36^\circ}$ (el suplementario mide $144^\circ=4\cdot36^\circ$).

\ex{3} Los ángulos alternos internos son iguales: $3x+10=70\ \Rightarrow\ 3x=60\ \Rightarrow\ x=\res{20^\circ}$.


% =====================================================
\seccio{6}{Polígonos: elementos y clasificación}

\ex{1} \begin{tasks}(3)
  \task \rest{pentágono}
  \task \rest{octógono}
  \task \rest{hexágono}
  \task \rest{decágono}
  \task \rest{heptágono}
  \task \rest{dodecágono}
\end{tasks}

\ex{1} \begin{tasks}(2)
  \task 5 lados, \rest{convexo} y \rest{regular}
  \task 5 lados, \rest{cóncavo} e \rest{irregular}
  \task 4 lados, \rest{convexo} e \rest{irregular} (rectángulo)
  \task 4 lados, \rest{convexo} e \rest{irregular}
\end{tasks}

\ex{1} De cada vértice salen $n-3$ diagonales; en total, $\dfrac{n\cdot(n-3)}{2}$.\begin{tasks}(4)
  \task 1 y \res{2}
  \task 2 y \res{5}
  \task 5 y \res{20}
  \task 7 y \res{35}
\end{tasks}

\ex{2} \begin{tasks}(1)
  \task \rest{Cierto} (y también todos los ángulos iguales).
  \task \rest{Falso}: sus ángulos no son todos iguales.
  \task \rest{Cierto}.
  \task \rest{Falso}: no tiene ninguna ($\frac{3\cdot0}{2}=0$).
\end{tasks}

\ex{2} Buscamos $n$ con $n\cdot(n-3)=28$: \ $7\cdot4=28$ \ $\Rightarrow$ \rest{heptágono}. \ Con $n\cdot(n-3)=54$: \ $9\cdot6=54$ \ $\Rightarrow$ \rest{eneágono}.

\ex{3} $n-3=9\ \Rightarrow\ n=\res{12}$ (dodecágono). \ En total: $\dfrac{12\cdot9}{2}=\res{54}$ diagonales.


% =====================================================
\seccio{7}{Triángulos}

\ex{1} \begin{tasks}(4)
  \task \rest{equilátero}
  \task \rest{isósceles}
  \task \rest{escaleno}
  \task \rest{isósceles}
\end{tasks}

\ex{1} \begin{tasks}(2)
  \task $180^\circ-100^\circ=\res{80^\circ}$, \rest{acutángulo}
  \task $180^\circ-125^\circ=\res{55^\circ}$, \rest{rectángulo}
  \task $180^\circ-145^\circ=\res{35^\circ}$, \rest{obtusángulo}
  \task $180^\circ-90^\circ=\res{90^\circ}$, \rest{rectángulo} (e isósceles)
\end{tasks}

\ex{1} El lado más largo tiene que ser menor que la suma de los otros dos.\begin{tasks}(4)
  \task $5<7$: \rest{sí}
  \task $6>5$: \rest{no}
  \task $9<10$: \rest{sí}
  \task $8=8$: \rest{no}
\end{tasks}

\ex{2} \begin{tasks}(1)
  \task $(180^\circ-50^\circ):2=\res{65^\circ}$ cada uno.
  \task $180^\circ-2\cdot72^\circ=\res{36^\circ}$.
  \task $(180^\circ-90^\circ):2=\res{45^\circ}$ cada uno.
\end{tasks}

\ex{2} \begin{tasks}(2)
  \task \rest{Sí}: ángulos de $90^\circ$, $45^\circ$ y $45^\circ$.
  \task \rest{No}: los ángulos de un equilátero miden $60^\circ$.
  \task \rest{Sí}: por ejemplo, $50^\circ$, $60^\circ$ y $70^\circ$.
  \task \rest{No}: los ángulos de un equilátero miden $60^\circ$.
\end{tasks}

\ex{2} Los dos ángulos agudos suman $90^\circ$: \ $89^\circ\,60'-27^\circ\,35'=\res{62^\circ\,25'}$.

\ex{2} \begin{tasks}(1)
  \task El \rest{circuncentro} (corte de las mediatrices).
  \task El \rest{ortocentro}.
  \task El \rest{baricentro} (corte de las medianas).
  \task El \rest{incentro} (corte de las bisectrices).
\end{tasks}

\ex{3} $x+(x+15)+(x+30)=180\ \Rightarrow\ 3x+45=180\ \Rightarrow\ x=45$. \ Ángulos: \rest{$45^\circ$, $60^\circ$ y $75^\circ$}: \rest{acutángulo} y escaleno.

\ex{3} El tercer lado debe ser menor que $5+12=17$ y mayor que $12-5=7$: \ \rest{8, 9, 10, …, 16\,cm} (9 posibilidades).

\ex{3} El ángulo interior de ese vértice mide $180^\circ-130^\circ=50^\circ$. \ El tercero: $180^\circ-60^\circ-50^\circ=70^\circ$. \ Ángulos: \rest{$50^\circ$, $60^\circ$ y $70^\circ$}.


% =====================================================
\seccio{8}{Cuadriláteros}

\ex{1} \begin{tasks}(3)
  \task \rest{cuadrado}
  \task \rest{rombo}
  \task \rest{trapecio isósceles}
  \task \rest{rectángulo}
  \task \rest{romboide}
  \task \rest{trapecio rectángulo}
\end{tasks}

\ex{1} \begin{tasks}(3)
  \task $\res{80^\circ}$
  \task $\res{105^\circ}$
  \task $\res{85^\circ}$
\end{tasks}

\ex{1} \begin{tasks}(2)
  \task \rest{cuadrado y rectángulo}
  \task \rest{cuadrado y rombo}
  \task \rest{cuadrado y rectángulo}
  \task \rest{cuadrado y rombo}
\end{tasks}

\ex{2} \begin{tasks}(1)
  \task $\res{54^\circ,\ 126^\circ,\ 54^\circ,\ 126^\circ}$ (ángulos opuestos iguales; consecutivos suplementarios)
  \task $\res{112^\circ,\ 68^\circ,\ 112^\circ,\ 68^\circ}$
  \task $\res{90^\circ,\ 90^\circ,\ 63^\circ,\ 117^\circ}$ ($180^\circ-63^\circ=117^\circ$)
\end{tasks}

\ex{2} \begin{tasks}(1)
  \task \rest{Cierto}: tiene los cuatro ángulos rectos.
  \task \rest{Falso}: el rombo no tiene por qué tener ángulos rectos.
  \task \rest{Cierto}.
  \task \rest{Falso}: si tiene tres, el cuarto también mide $90^\circ$ y sería un rectángulo.
\end{tasks}

\ex{3} \begin{tasks}(1)
  \task $10x=360\ \Rightarrow\ x=36$: \ \rest{$36^\circ$, $72^\circ$, $108^\circ$ y $144^\circ$}.
  \task Dos ángulos consecutivos suman $180^\circ$: $x+3x=180\Rightarrow x=45$. \ \rest{$45^\circ$, $135^\circ$, $45^\circ$ y $135^\circ$}.
\end{tasks}


% =====================================================
\seccio{9}{Ángulos de los polígonos. Polígonos regulares}

\ex{1} \begin{tasks}(5)
  \task $\res{540^\circ}$
  \task $\res{720^\circ}$
  \task $\res{1080^\circ}$
  \task $\res{1440^\circ}$
  \task $\res{1800^\circ}$
\end{tasks}

\ex{1} \begin{tasks}(3)
  \task $\res{60^\circ}$ y $\res{120^\circ}$
  \task $\res{90^\circ}$ y $\res{90^\circ}$
  \task $\res{108^\circ}$ y $\res{72^\circ}$
  \task $\res{120^\circ}$ y $\res{60^\circ}$
  \task $\res{135^\circ}$ y $\res{45^\circ}$
  \task $\res{144^\circ}$ y $\res{36^\circ}$
\end{tasks}

\ex{2} $n=360:24=\res{15}$ lados. \ Ángulo interior: $180^\circ-24^\circ=\res{156^\circ}$ \ (o bien $\frac{13\cdot180^\circ}{15}=156^\circ$).

\ex{2} Ángulo central: $180^\circ-160^\circ=20^\circ$; \ $n=360:20=\res{18}$ lados.

\ex{2} $(n-2)\cdot180=2340\ \Rightarrow\ n-2=13\ \Rightarrow\ n=\res{15}$ lados.

\ex{3} Con $100^\circ$, el ángulo central mediría $80^\circ$ y $360:80=4,5$ no es un número natural: \rest{no}. \ Con $150^\circ$: $360:30=12$: \rest{sí}, el dodecágono regular.

\ex{3} Los ángulos que se encuentran en un vértice tienen que sumar $360^\circ$: \rest{triángulos} ($6\cdot60^\circ$), \rest{cuadrados} ($4\cdot90^\circ$) y \rest{hexágonos} ($3\cdot120^\circ$). Con pentágonos no ($3\cdot108^\circ=324^\circ$) ni con octógonos solos. Combinándolos, \rest{sí}: dos octógonos y un cuadrado, $135^\circ+135^\circ+90^\circ=360^\circ$.

\ex{3} Cada ángulo central mide $360^\circ:6=60^\circ$. Los dos radios son iguales, así que los otros dos ángulos son iguales: $(180^\circ-60^\circ):2=60^\circ$. Los seis triángulos son \rest{equiláteros}.


\end{document}
